\documentclass[11pt,a4paper]{article}
\usepackage{amsmath,amssymb,amsthm}
\usepackage[margin=2.5cm]{geometry}
\usepackage{hyperref}

\newtheorem{definition}{Definition}[section]
\newtheorem{theorem}[definition]{Theorem}
\newtheorem{proposition}[definition]{Proposition}
\newtheorem{lemma}[definition]{Lemma}
\newtheorem{corollary}[definition]{Corollary}
\newtheorem{conjecture}[definition]{Conjecture}
\newtheorem{remark}[definition]{Remark}

\title{Hyperseed Formalization:\\
  The Folly of Statistically Watermarking AI-Generated Text}
\author{ProtomegaTron (automated formalization)}
\date{2026-09-17}

\begin{document}
\maketitle

\begin{abstract}
We formalize the intellectual structure of Ben Goertzel's Substack article
``The Folly of Statistically Watermarking AI-Generated Text'' (2026-07-21)
within the Hyperseed ontology. The article argues that statistical text
watermarking (SynthID-style) cannot achieve what its proponents claim and
will instead produce a surveillance infrastructure concentrating control
over text production. Each structural insight maps onto Hyperseed structures:
statistical watermarking as 0-cell observable property (vs.\ cryptographic
provenance as directed history), the arms race as a centralizing cocycle
trending toward frozen directed type, the split ecosystem as fiber bundle
fracture, reasoning-path watermarks as higher-cell fingerprinting, and
the detector as cocycle controller.
\end{abstract}

\tableofcontents

%% =================================================================
\section{Statistical watermarking as 0-cell property}
\label{sec:watermark}
%% =================================================================

\begin{definition}[Dice-loading, not signing --- claims 1--5, 26]
\label{def:watermark}
SynthID-Text works by biasing the model's word choices during generation
toward patterns detectable by a companion algorithm. This is a
\emph{fingerprint on one generation process}, not a signature on the
content. In Hyperseed terms, the watermark is a 0-cell observable
property:
\begin{itemize}
\item It examines observable properties of the current text state
  (word-choice distributions) without accessing the directed history
  of how the text was produced.
\item It tells you about one particular generation path --- not about
  the full causal history of ideas, research, reasoning, and authorship.
\item Constrained writing (proofreading, factual reporting) carries
  almost no watermark because the model has no distributional freedom
  to exploit.
\item Short passages lack sufficient statistical evidence for confident
  detection.
\end{itemize}
This is the same detection-vs-provenance distinction from the Provenance
Layer article: detection examines 0-cells (current observable state);
provenance asks about 1-cells and higher (transformation history).
\end{definition}

%% =================================================================
\section{Five removal pathways: increasing cell level}
\label{sec:removal}
%% =================================================================

\begin{proposition}[Removal taxonomy --- claims 6--11]
\label{prop:removal}
The five removal pathways form a hierarchy by the cell level at which
the watermark operates:
\begin{enumerate}
\item \textbf{Closed API + paraphrase} (0-cell removal): the watermark
  lives in surface word choices; replacing enough word choices destroys
  it. No key knowledge needed --- the attacker just operates at the
  0-cell level.
\item \textbf{Open weights + sampler swap} (0-cell removal, trivial):
  whoever controls inference controls everything downstream. The
  watermark lives in the sampling procedure; swapping the sampler
  removes it entirely.
\item \textbf{Weight-baked + retraining} (0.5-cell removal): the
  watermark is baked into the model's output distribution, surviving
  sampler swap. Capability-preserving retraining on clean data removes
  it. Harder but very plausible for a motivated actor.
\item \textbf{Reasoning-path watermarks} (1-cell removal): the watermark
  encodes preferences in reasoning strategy, explanation structure,
  problem decomposition --- 1-cells in the cognitive directed type.
  Removing this requires changing how the model thinks, not just how
  it talks. This is \emph{higher-cell fingerprinting}: harder to remove
  because it touches the directed structure rather than just the
  observable surface.
\item \textbf{Detector feedback} (adaptive): with detector access,
  removal becomes adaptive optimization --- preserve capabilities while
  driving the detector score down. This is the arms-race endpoint
  where all previous categories merge into an optimization problem.
\end{enumerate}
\end{proposition}

\begin{remark}[Watermark radioactivity]
A student model trained on a watermarked teacher inherits faint watermark
traces (``radioactivity'') because the student absorbs the teacher's
subtle lexical preferences. But the radioactivity decays with fresh data,
RL, or preference tuning --- these interventions rewrite the inherited
biases. The radioactivity is a 0-cell echo, not a robust signal.
\end{remark}

%% =================================================================
\section{The arms race as centralizing cocycle}
\label{sec:arms}
%% =================================================================

\begin{theorem}[Centralizing cocycle --- claims 12, 27]
\label{thm:arms}
The watermark arms race is an enduring co-evolutionary cycle:
\[
  \text{watermark}_n \to \text{scrubber}_n \to \text{watermark}_{n+1}
  \to \text{scrubber}_{n+1} \to \cdots
\]
Each turn of the ratchet increases the intrusiveness of the watermarking
scheme:
\[
  \text{sampling bias} \to \text{weight baking} \to \text{detector secrecy}
  \to \text{legal restrictions} \to \text{trusted-computing lockdown}
\]
In Hyperseed terms, this is a \emph{centralizing cocycle}: each transition
function $g_n$ in the arms race takes the system further from distributed
text production and closer to centralized control. The cocycle trends
toward frozen directed type --- the same pathology as in Seven Flavors,
where capturing the development environment prevents cooperative dynamics
and leaves only adversarial ones.

The arms race does not converge to a solution. It converges to a control
infrastructure.
\end{theorem}

%% =================================================================
\section{Cryptographic provenance as directed history}
\label{sec:provenance}
%% =================================================================

\begin{definition}[Accent vs.\ signature --- claims 13--16, 30]
\label{def:provenance}
Two fundamentally different mechanisms:
\begin{itemize}
\item \textbf{Statistical watermark} (accent): a detectable pattern in
  the surface text, operating at the 0-cell level. Cannot be deliberately
  verified or explicitly checked by the user. Fragile under transformation.
\item \textbf{Cryptographic credential} (signature): a deliberate,
  checkable declaration recording the transformation history of the
  artifact. Operates at the 1-cell level (who did what, in what order,
  with what tools). Tamper-evident. Standards like C2PA and protocols
  like OpenWater.
\end{itemize}

Signed provenance gives strong positive evidence when the chain is intact.
But the absence of a credential proves nothing about human origin ---
copy the text into a fresh document and the metadata is gone. Cryptography
authenticates a declared chain of custody; it cannot force every causal
influence on an idea to stay attached to it forever.

The fundamental distinction: statistical watermarking examines the
observable surface (0-cells); cryptographic provenance records the
directed history (1-cells and higher). This is the same
detection-vs-provenance distinction, now applied to text rather than media.
\end{definition}

%% =================================================================
\section{The split ecosystem as fiber bundle fracture}
\label{sec:split}
%% =================================================================

\begin{proposition}[Compliance perimeter --- claims 17--20, 28]
\label{prop:split}
The likely outcome of widespread watermarking deployment is a
\emph{compliance perimeter} splitting the text production fiber into
two disconnected components:
\begin{enumerate}
\item \textbf{Compliant side}: certified models, signed runtimes,
  approved enterprise services, watermarked outputs.
\item \textbf{Non-compliant side}: open models, modified derivatives,
  locally run systems, models deliberately cleaned to resist detection.
\end{enumerate}

This perimeter does not track the AI-human boundary. It tracks the
institutional-control boundary --- which text came through an approved
pipeline and which didn't.

In Hyperseed terms, this is a \emph{fiber bundle fracture}: the
transition functions between the two components are broken. Text from
one side cannot be authenticated by the other's standards. The result
is a class structure in text production --- approved-source text treated
as trustworthy by default; unapproved-source text treated as suspect
regardless of quality.
\end{proposition}

%% =================================================================
\section{Absence as guilt and the detector as chokepoint}
\label{sec:political}
%% =================================================================

\begin{proposition}[Absence $\neq$ guilt --- claims 15, 21, 29]
\label{prop:absence}
Two rhetorical sleights of hand:
\begin{enumerate}
\item ``This signed artifact came through an approved pipeline''
  $\Rightarrow$ ``this unsigned artifact must be suspect.''
\item ``This AI system is institutionally approved''
  $\Rightarrow$ ``its outputs are epistemically superior.''
\end{enumerate}
Both are instances of the absence-as-guilt fallacy --- the same design
refusal that OpenWater explicitly encodes: ``no absence-as-guilt.''
The absence of a watermark is absence of evidence, not evidence of AI
generation.
\end{proposition}

\begin{proposition}[Detector as cocycle controller --- claims 23, 32]
\label{prop:detector}
If watermark detection becomes legally mandated, the detector becomes a
chokepoint: whoever controls the detector controls what counts as
``AI-generated.'' In Hyperseed terms, the detector is the
\emph{cocycle controller} --- it controls the transition function between
the two halves of the split ecosystem:
\[
  g_{\text{detector}} : E_{\text{compliant}} \to E_{\text{non-compliant}}
\]
This is the same chokepoint-capture dynamic as centralized provenance
infrastructure (the failure mode OpenWater's decentralized resolution
is designed to prevent).
\end{proposition}

%% =================================================================
\section{Cross-references}
%% =================================================================

\begin{remark}[Connections to other formalizations]
\begin{itemize}
\item \textbf{Provenance Layer} (2026-08-25): detection vs.\ provenance
  as observables vs.\ directed history. Here the same distinction applied
  to text: statistical watermarking is detection (0-cells); cryptographic
  provenance is directed history (1-cells).
\item \textbf{Proof of Humanity} (2026-08-26): OpenWater's four design
  refusals, including ``no absence-as-guilt.'' Here that refusal is
  violated by the assumption that unwatermarked text is suspect.
\item \textbf{Seven Flavors} (2026-07-01): frozen directed type as evil
  mode. The watermark arms race converges toward frozen directed type ---
  centralized control over text production that eliminates cooperative
  dynamics.
\item \textbf{Sanders Ban} (2026-09-05): state capture of AI development.
  Mandatory watermark detection is the text-production version of the
  Sanders ban --- using compliance infrastructure to control who gets to
  produce what.
\item \textbf{Big Tech's Bets} (2026-08-02): too-strategic-to-fail
  dynamics. If watermarking infrastructure becomes too entrenched to
  dismantle, it becomes a permanent control layer regardless of whether
  it achieves its stated goals.
\item \textbf{Zuck's Future} (2026-08-10): access vs.\ ownership. The
  compliance perimeter creates distributed access to approved AI text
  without distributed power over what counts as approved.
\item \textbf{$(f,c)$-lossy critique} (LTM 5a4d150b): binary
  watermark-detected/not-detected is the ultimate scalar collapse ---
  reducing the rich provenance question to a single bit.
\end{itemize}
\end{remark}

\begin{conjecture}[Reasoning-path watermark persistence]
Conjecture: reasoning-path watermarks (1-cell fingerprints) are
exponentially harder to remove than surface watermarks (0-cell fingerprints),
but also exponentially harder to \emph{inject} reliably. The difficulty
of injection limits their practical deployment, creating an asymmetry:
the most robust watermark type is also the least deployable, while the
most deployable type (sampling-layer) is the least robust. This
asymmetry means the watermarking approach is fundamentally
self-undermining at every cell level.
\end{conjecture}

\end{document}
